Multi-agent large language model systems often rely on one model to critique
another, but a revised answer alone does not reveal how the receiving model
internally processes that feedback. We study feedback uptake in a
Solver--Critic system using MuSiQue, a multi-hop question-answering benchmark
with known answers. A Gemma 3 4B instruction-tuned Solver first answers each
question, a Qwen3 4B instruction-tuned Critic independently reviews that
answer, and the Solver then chooses between its original response and the
Critic's recommendation. We capture Solver residual-stream activations across
nine transformer layers and use sparse autoencoders (SAEs) to represent those
activations as sparse features. Linear probes identify features associated
with adoption versus retention, while held-out causal interventions test
whether altering selected features changes the Solver's decision. The
collected candidate-presentation dataset contains 22,355 sampled questions;
the current labeling policy identifies 5,834 adoptions, 4,810 retentions,
8,451 agreements, and 3,260 episodes excluded from the primary behavioral
analysis. This design separates predictive association from causal evidence
and provides a reproducible framework for investigating when language models
defer to, resist, or agree with another model's feedback.
